\chapter{Smart Order Routing}\label{ch:mx:smart-order-routing}

A router splits an order across five venues. Sent at once, the pieces arrive microseconds apart, and quotes on the later venues vanish in between:
the first fill told fast traders what was coming. This chapter builds a router on the multi-venue simulator, measures what the order loses when its
pieces arrive one after another and what it keeps when they are timed to arrive together, ranks venues by what they deliver rather than what they
show, allocates a passive order across venues by Cont and Kukanov's programme, and ends with the fees and the rules that protect an order from being
read.

\section{What a router decides}

\begin{definition}[Smart order router]\label{def:mx:smart-order-routing:sor}
A \emph{smart order router}\index{smart order router} splits a \omterm{def:mx:the-almgren-chriss-framework:parent}{child order} across venues and decides, for each piece, the venue, the price, the size,
the order type and the moment it is sent, from its view of the venues' quotes, fees, latencies and past fills.
\end{definition}

Chapter~17 decided whether to rest or cross; the router decides where. Its inputs are the venues' books (chapter~8's fragmented market, \omterm{def:mx:fragmentation-and-routing:protected}{protected
quotes} and the order protection rule of One Quant Book 1, chapter~9), their fees (chapter~7's \omterm{def:mx:fees-rebates-and-venue-economics:adjusted}{fee-adjusted price}), its own latency to each venue
(One Quant Book 7, chapter~18's order-entry latency) and statistics of what past orders got. Foucault and Menkveld (2008) studied the entry of the
London Stock Exchange against Euronext in Dutch stocks: the consolidated book was deeper after entry, and where more trades went through the entrant's
better quotes, it supplied less liquidity, which is why routers that actually reach every \omterm{def:mx:fragmentation-and-routing:protected}{protected quote} matter.

The experiment of this chapter uses five venues, A to E, in \texttt{firm.exchsim}. The router sits next to A and reaches the venues in 50, 120, 200,
280 and 360 microseconds. Each venue shows 500 shares at the ask of 100.01: 200 from slow providers and 300 from one fast provider, whose machine sits
at A and reaches the other venues over faster links (45, 95, 145 and 195 microseconds). Behind them each venue has 1\,000 shares at 100.02. The
fast provider follows one rule: when one of its orders fills, it cancels its orders on every other venue. Take fees differ: 0.30 cents a share at A,
B and D, 0.10 at C, and a rebate of 0.10 to the taker at E (an inverted venue, chapter~7).

\section{Venue ranking and fill probabilities}

\begin{definition}[Venue ranking]\label{def:mx:smart-order-routing:ranking}
A \emph{venue ranking}\index{venue ranking} orders the venues by the expected cost of routing a share there, from the fill probability the router has
measured on its own orders, the fee, and the cost of what does not fill.
\end{definition}

A router that believes displayed quotes ranks venues by \omterm{def:mx:fees-rebates-and-venue-economics:adjusted}{fee-adjusted price}. One that keeps statistics learns what the display is worth. After 100
sprays (below) with a latency jitter of 10\%, the router's own orders filled 100\% of the shares sent to A and 44\%, 41\%, 40\% and 40\% of those sent
to B, C, D and E. With an unfilled share costing a tick later, the expected cost of a share sent to each venue ranks A first (0.30 cents), then E (0.56),
C (0.63), B (0.69) and D (0.72): the displayed rebate of E is worth less than the certainty of A. Timed to arrive together, the same orders filled
everywhere, and the ranking reverts to the fees: E ($-0.10$), C (0.10), then A, B and D (0.30). A \omterm{def:mx:smart-order-routing:ranking}{venue ranking} is a statement about the router's
own routing as much as about the venues.

\omcode{code/firm/sor/firm_sor.py}{108}{121}{The router's sends: each leg of the sweep is scheduled at the start time plus a delay that is zero for a spray and, for a synchronised sweep, the gap between the venue's latency and the farthest one's.}\label{lst:mx:smart-order-routing:router}

\section{Aggressive routing and quote fade}

\begin{definition}[Spray order, quote fade]\label{def:mx:smart-order-routing:spray}
A \emph{spray order}\index{spray order} sends marketable pieces to several venues at the same moment, so that they arrive in the order of the router's
latencies. \emph{Quote fade}\index{quote fade} is the disappearance of displayed quotes between the moment a router sees them and the moment its
orders arrive, typically because a liquidity provider cancels on seeing a fill elsewhere.
\end{definition}

\begin{definition}[Synchronised routing]\label{def:mx:smart-order-routing:sync}
\emph{Synchronised routing}\index{synchronised routing} delays each piece of a multi-venue order by the difference between its venue's latency and the
largest, so that all pieces arrive at their venues at about the same time and no fill can warn anyone before the last piece lands.
\end{definition}

\Cref{fig:mx:smart-order-routing:race} is the race. The spray reaches A at 50 microseconds and fills there; the fast provider learns of its fill 5
microseconds later and its cancels reach B at 100, C at 150, D at 200 and E at 250, each before the router's piece. The router gets A's 500 shares and
the slow 200 on each other venue: 1\,300 of 2\,500, a fill ratio of 52\%. The 1\,200 shares left go to the next level in a clean-up sweep, and the whole
order costs 0.73 cents a share above the displayed ask, fees included. Synchronised, every piece arrives at 360 microseconds: the fill ratio is 100\%
and the cost 0.18 cents, the average take fee.

\begin{omfigure}
\begin{tikzpicture}
\begin{axis}[width=10cm,height=5cm,xlabel={\small microseconds after the router's decision},xmin=0,xmax=400,ymin=0.4,ymax=5.6,
  ytick={1,2,3,4,5},yticklabels={E,D,C,B,A},tick label style={font=\footnotesize},grid=major,grid style={gray!25},
  legend columns=3,legend style={font=\footnotesize,at={(0.5,-0.32)},anchor=north,draw=none,fill=none,
  /tikz/every even column/.append style={column sep=8pt}},table/col sep=comma]
\addplot[omDef,only marks,mark=*,mark size=2.4pt] table[x=spray,y=i]{figdata/microstructure/18-smart-order-routing/race.csv};
\addplot[omProp,only marks,mark=x,mark size=4pt,very thick] table[x=cancel,y=i]{figdata/microstructure/18-smart-order-routing/race.csv};
\addplot[omThm,only marks,mark=square*,mark size=2.2pt] table[x=sync,y=i]{figdata/microstructure/18-smart-order-routing/race.csv};
\legend{spray arrives,fast provider's cancel arrives,synchronised arrival}
\end{axis}
\end{tikzpicture}

\omcaption{The race at each venue: when the router's sprayed piece arrives, when the fast provider's cancel arrives after its fill at A, and when a
synchronised piece arrives. Every cancel beats the spray (at E the sprayed and synchronised pieces arrive together); none can beat the synchronised pieces. Data: \texttt{mx\_sor.race}, from the latencies of
the experiment.}\label{fig:mx:smart-order-routing:race}
\end{omfigure}

\omcode{code/firm/sor/firm_sor.py}{164}{170}{The fast provider's rule: on the first fill of one of its orders, cancel every other order it has on the other venues.}\label{lst:mx:smart-order-routing:fader}

Synchronisation works only as well as the router knows its latencies. \Cref{fig:mx:smart-order-routing:jitter} adds a random factor to each piece's
latency (lognormal, mean one), 100 trials per level. Up to 10\% the synchronised sweep still fills everything; at 20\% it fills 99.3\% of the shares on
average and the whole level on 94\% of trials; at 40\%, 90.3\% and 37\%. The spray gains from the noise, since it reshuffles the arrivals, but only to
68.6\%. Royal Bank of Canada patented the method, in claims that describe timing parameters determined from latencies ``so as to cause synchronized
arrival'' of the pieces at the venues (US patent 8\,489\,747, granted in 2013); its description names the problem as orders arriving at faster
exchanges before slower ones, where others can detect them.

\begin{omfigure}
\begin{tikzpicture}
\begin{axis}[width=9.4cm,height=5.4cm,xlabel={\small latency jitter (lognormal $\sigma$)},ylabel={\small fill ratio at the displayed ask (\%)},
  tick label style={font=\footnotesize},grid=major,grid style={gray!25},xmin=-0.02,xmax=0.42,ymin=40,ymax=105,
  xtick={0,0.05,0.1,0.2,0.4},xticklabels={0,0.05,0.1,0.2,0.4},
  legend columns=2,legend style={font=\footnotesize,at={(0.5,-0.3)},anchor=north,draw=none,fill=none,
  /tikz/every even column/.append style={column sep=8pt}},table/col sep=comma]
\addplot[omProp,very thick,mark=*] table[x=jitter,y=spray]{figdata/microstructure/18-smart-order-routing/jitter.csv};
\addplot[omDef,very thick,mark=square*] table[x=jitter,y=sync]{figdata/microstructure/18-smart-order-routing/jitter.csv};
\legend{spray,synchronised}
\end{axis}
\end{tikzpicture}

\omcaption{Fill ratio of a sweep of the five venues' 2\,500 displayed shares against the jitter of the router's latencies: 100 trials per level in the
simulated market with a fading fast provider. Data: \texttt{mx\_sor.sweep\_study}.}\label{fig:mx:smart-order-routing:jitter}
\end{omfigure}

An order smaller than the displayed liquidity makes the router choose venues. For 1\,000 shares (with 10\% jitter), a router that picks the cheapest
fees goes to E and C, the two farthest venues, and loses almost nothing even when it sprays (99.8\%, 0.004 cents a share), because its first fill lands
where the fast provider's cancels are slow to reach. A router that picks the nearest venues goes to A and B and, spraying, fills 71.2\% at 0.59 cents;
synchronised, 100\% at 0.30. Which piece arrives first matters more than which venue is cheapest.

\section{Passive allocation across venues}

A passive order faces the opposite problem: where to rest so as to be filled. Cont and Kukanov (2017) wrote it as a convex programme. With $X$ shares
to buy, a \omterm{def:mx:the-limit-order-book:orders}{market order} $M$ and \omterm{def:mx:the-limit-order-book:orders}{limit orders} $L_k$ behind queues $Q_k$, and $\xi_k$ the shares that will leave the front of queue $k$ during the horizon
(\omterm{def:mx:the-limit-order-book:orders}{market orders} and cancellations, random), the \omterm{def:mx:the-limit-order-book:orders}{limit order} at $k$ fills $\min(L_k,(\xi_k-Q_k)^+)$ and the cost against the mid is
\[
(h+f)M-\sum_k(h+r_k)\min\bigl(L_k,(\xi_k-Q_k)^+\bigr)+\lambda_u(X-A)^++\lambda_o(A-X)^+,
\]
with $h$ the half-spread, $f$ the take fee, $r_k$ the rebate, $A$ the shares executed, and penalties $\lambda_u$, $\lambda_o$ for executing too few or
too many. Its expectation is convex when the penalties exceed the spread and fees; \texttt{firm\_sor.ck\_allocate} minimises the sample average over
simulated outflows.

\omcode{code/firm/sor/firm_sor.py}{178}{186}{The Cont and Kukanov cost of an allocation on a set of outflow scenarios: each limit order fills what leaves its queue beyond the orders ahead of it, and penalties price the shares executed too few or too many.}\label{lst:mx:smart-order-routing:ck}

Take four lit venues and a dark one for a 1\,000-share buy over a minute: queues of 1\,500, 2\,500, 300 and 800 shares at the lit bids and none in the
dark pool; outflows with means of 2\,000, 2\,200, 900, 1\,000 and 400 shares (lognormal, with a common factor); rebates of 0.20, 0.25, $-0.05$ and 0.20
cents (the third venue is inverted); the dark venue fills at the mid, half a spread (0.5 cents) worse than a fill at the bid, without fee; a take fee of
0.30 cents, and penalties of 1.2 cents a share either way. The optimum is fitted on 400 scenarios and every rule is scored on 20\,000 fresh ones
(\cref{tab:mx:smart-order-routing:passive}).

\begin{table}[ht]
\centering
\small
\begin{tabular}{@{}lrrl@{}}
\toprule
rule & cost (cents a share) & filled & allocation (market; A, B, C, D, dark) \\
\midrule
\omterm{def:mx:the-limit-order-book:orders}{market order} & 0.80 & 100.0\% & 1\,000; 0, 0, 0, 0, 0 \\
highest rebate & 0.77 & 22.0\% & 0; 0, 1\,000, 0, 0, 0 \\
shortest queue & 0.74 & 38.7\% & 0; 0, 0, 0, 0, 1\,000 \\
proportional to room & 0.20 & 61.3\% & 0; 266, 171, 255, 140, 169 \\
Cont--Kukanov & 0.03 & 76.4\% & 0; 405, 373, 406, 401, 193 \\
Cont--Kukanov, no dark venue & 0.06 & 69.7\% & 0; 544, 478, 522, 478, 0 \\
\bottomrule
\end{tabular}
\caption{A 1\,000-share passive buy over five venues: expected cost against the mid including penalties, share of the order filled, and the
allocation, scored out of sample. Data: \texttt{mx\_sor.passive\_study}.}\label{tab:mx:smart-order-routing:passive}
\end{table}

The optimum posts 1\,778 shares for 1\,000 wanted: it overbooks, because each venue fills only part of its order and the penalty for too many is
paid only on the scenarios where most venues fill. It sends nothing at the market. Chasing the highest rebate puts the whole order behind the longest
queue and fills 22\%; Battalio, Corwin and Jennings (2016) found in real data that \omterm{def:mx:the-limit-order-book:orders}{limit orders} routed to venues paying the largest rebates execute
worse, and that brokers who route for rebates cannot have it all. The dark venue is worth 0.03 cents a share here: its fills are worse than a
passive fill at the bid but come without a queue. Maglaras, Moallemi and Zheng (2021) showed that when many traders route this way, the queues of the
venues move together, and the market behaves like one aggregate queue.

\section{Fee-aware routing and anti-gaming}

A fee-aware router compares prices net of fees (chapter~7): at equal displayed prices it takes the inverted venue first and rests on the venue with
the best rebate. The two lessons above temper both: the cheapest venue to take from is worth nothing if its quote has faded by the time the order
arrives, and the best rebate is worth nothing behind a queue that will not clear.

\begin{definition}[Anti-gaming logic]\label{def:mx:smart-order-routing:antigaming}
\emph{Anti-gaming logic}\index{anti-gaming logic} is the part of a router that makes its orders hard to detect and exploit: synchronised arrivals,
minimum fill quantities in dark venues, randomised child sizes and timing, avoiding venues whose fills are followed by adverse moves, and not routing
to venues that let other participants learn about resting orders.
\end{definition}

Chapter~9 measured how \omterm{def:mx:dark-pools-and-blocks:ping}{pinging orders} find a large dark order and what the leak costs. A router answers with a dark-first leg that carries a minimum
fill quantity, so that a 100-share ping cannot trade with it, and with child sizes drawn around their target instead of repeated
(\texttt{firm\_sor.dark\_first}, \texttt{child\_sizes}); it answers \omterm{def:mx:smart-order-routing:spray}{quote fade} with synchronised arrivals; and it keeps, per venue, the markouts of its
fills (One Quant Book 7, chapter~23) to stop routing to venues where they are consistently bad. Latency arbitrage, the business of the fast provider
of this chapter, is One Quant Book 11's subject.

\section{Tutorial: a router on five venues}\label{tut:mx:smart-order-routing:tut}

\begin{tutorial}
\textbf{Goal.} Sweep five simulated venues against a fading fast provider, with and without arrival-time alignment; rank the venues; allocate a
passive order by Cont and Kukanov's programme.
\textbf{End state:} \cref{fig:mx:smart-order-routing:race,fig:mx:smart-order-routing:jitter}, \cref{tab:mx:smart-order-routing:passive} and the numbers
of sections 2 to 4.
\begin{tutsteps}
\item \textbf{Venues.} \texttt{mx\_sor.VENUES}, \texttt{race()}; \texttt{firm\_sor.Venue}, \texttt{send\_delays}.
\item \textbf{Sweep.} \texttt{trial(sync, jitter, seed)} with \texttt{firm\_sor.Router} and \texttt{Fader}; \texttt{sweep\_study()}.
\item \textbf{Ranking.} \texttt{VenueStats.add}, \texttt{fill\_ratio}, \texttt{rank}; \texttt{partial\_study()}.
\item \textbf{Passive.} \texttt{outflow(n, seed)}, \texttt{ck\_cost}, \texttt{ck\_allocate}; \texttt{passive\_study()}; draw with \texttt{fig\_sor.py}.
\end{tutsteps}
\textbf{What to change next.} Give the fast provider a second machine at C; let the router learn its latencies from acknowledgements; run the passive
allocation on simulated venues instead of scenarios.
\end{tutorial}

\section{Build: smart order router}\label{bld:mx:smart-order-routing:sor}

\begin{build}
\bfield{Purpose} The routing layer of chapter~28's \omterm{def:mx:benchmark-algorithms:algo}{execution algorithm} and of the desk of chapter~20: sweeps, passive allocation and venue statistics.
\bfield{Interface} \texttt{Venue(name, entry\_ns, take, make, dark)}, \texttt{send\_delays(venues, sync, margin\_ns)}, \texttt{sweep(quotes, qty,
fees)}, \texttt{VenueStats} (\texttt{add}, \texttt{fill\_ratio}, \texttt{rank}), \texttt{Router(legs, venues, t0\_ns, sync, cleanup)},
\texttt{Fader(price, qty, venues, start\_ns)}, \texttt{ck\_cost}, \texttt{ck\_allocate}, \texttt{child\_sizes(qty, n, rng, spread)},
\texttt{dark\_first(qty, dark, min\_qty)}.
\bfield{Rules} Buys (sells mirror them); fees per share, positive when paid; client order identifiers unique across venues (the simulator numbers them
per session); a clean-up sweep one tick higher after every leg has reported.
\bfield{Acceptance tests} \texttt{code/firm/sor/tests/}: delays and fee ordering; venue statistics and ranking; on two simulated venues the fast
provider's \omterm{def:mx:smart-order-routing:spray}{quote fades} before a spray and not before a synchronised sweep, and the clean-up completes the order; the Cont--Kukanov cost on a hand
example and an optimum no worse than simple allocations; anti-gaming sizes.
\bfield{Stretch} Learned latencies; a routing table by time of day; per-venue markouts in the ranking.
\end{build}

\begin{omsources}
\item T.~Foucault and A.~J.~Menkveld, ``Competition for order flow and smart order routing systems'', \emph{Journal of Finance} 63(1), 2008.
\item Royal Bank of Canada, ``Synchronized processing of data by networked computing resources'', US patent 8\,489\,747 B2, 2013.
\item R.~Battalio, S.~A.~Corwin and R.~Jennings, ``Can brokers have it all? On the relation between make-take fees and limit order execution
quality'', \emph{Journal of Finance} 71(5), 2016.
\item R.~Cont and A.~Kukanov, ``Optimal order placement in limit order markets'', \emph{Quantitative Finance} 17(1), 2017.
\item C.~Maglaras, C.~C.~Moallemi and H.~Zheng, ``Queueing dynamics and state space collapse in fragmented limit order book markets'',
\emph{Operations Research} 69(4), 2021.
\end{omsources}

\section{Exercises}

\begin{exercise}[$\star$]\label{exo:mx:smart-order-routing:1}
With the experiment's latencies, how long must the router hold the piece for A in a synchronised sweep, and the piece for C?
\end{exercise}

\begin{exercise}[$\star$]\label{exo:mx:smart-order-routing:2}
The fast provider's machine moves to C. Its links reach A, B, D and E in 95, 50, 50 and 100 microseconds, and it hears of its fills at C in
5 microseconds. Which of its quotes does a spray from the router still find?
\end{exercise}

\begin{exercise}[$\star$]\label{exo:mx:smart-order-routing:3}
Compute the fee-adjusted cost of the spray: 1\,300 shares at 100.01 (500 at A, 200 at each other venue) and a clean-up of 1\,200 at 100.02 (1\,000
at A and 200 at B), with the fees of the experiment.
\end{exercise}

\begin{exercise}[$\star\star$]\label{exo:mx:smart-order-routing:4}
Why does the spray's fill ratio rise with latency jitter?
\end{exercise}

\begin{exercise}[$\star\star$]\label{exo:mx:smart-order-routing:5}
Why does the Cont--Kukanov optimum post more shares than it wants, and what stops it from posting many more?
\end{exercise}

\begin{exercise}[$\star\star$]\label{exo:mx:smart-order-routing:6}
A router ranks venues by their displayed rebate. What does the passive study say it will get, and why?
\end{exercise}

\begin{exercise}[$\star\star\star$]\label{exo:mx:smart-order-routing:7}
\emph{Coding.} Rerun the passive study with an over-fill penalty of 0.5 cents instead of 1.2. What does the optimum post, and what does it cost out of
sample?
\end{exercise}

\begin{exercise}[$\star\star\star$]\label{exo:mx:smart-order-routing:8}
\emph{Find the flaw.} ``Venue E showed 500 shares and we got 200: E's quotes are fake, stop routing there.''
\end{exercise}

\section{Problem: The Order That Arrived After the Quotes Had Gone}

\begin{problem}[{Weekend problem --- the order that arrived after the quotes had gone}]\label{pb:mx:smart-order-routing:1}
A desk's sweeps fill half of what the consolidated book shows. It asks why, and what a better router would change.

\medskip\noindent\textbf{Part I --- The router's problem.}
\begin{enumerate}
\item Define a \omterm{def:mx:smart-order-routing:sor}{smart order router} and list what it decides.
\item Describe the five venues, their latencies and fees, and the fast provider.
\item Define \omterm{def:mx:smart-order-routing:spray}{quote fade} and a \omterm{def:mx:smart-order-routing:spray}{spray order}.
\item Why did Foucault and Menkveld's evidence make routing to every \omterm{def:mx:fragmentation-and-routing:protected}{protected quote} matter?
\end{enumerate}

\medskip\noindent\textbf{Part II --- The sweep.}
\begin{enumerate}[resume]
\item Draw the race of the spray at each venue.
\item What fill ratio does the spray get, and what does the whole order cost?
\item Define \omterm{def:mx:smart-order-routing:sync}{synchronised routing} and compute the delays.
\item \emph{State the named result}: the fill ratio of the sweep with and without arrival-time alignment, and its fee-adjusted cost.
\item How does latency jitter change both?
\end{enumerate}

\medskip\noindent\textbf{Part III --- Ranking and choosing venues.}
\begin{enumerate}[resume]
\item Give the router's measured fill ratios by venue after spraying.
\item Rank the venues with and without synchronisation, and explain the difference.
\item For 1\,000 shares, compare routing by fee and by distance, sprayed and synchronised.
\item Why did the fee-first spray lose almost nothing?
\end{enumerate}

\medskip\noindent\textbf{Part IV --- Passive allocation and the answer.}
\begin{enumerate}[resume]
\item Write the Cont--Kukanov cost.
\item Give the cost and fill of the \omterm{def:mx:the-limit-order-book:orders}{market order}, the highest rebate, the shortest queue, the proportional rule and the optimum.
\item What is the dark venue worth here?
\item What did Battalio, Corwin and Jennings find about rebates and execution quality?
\item Name three anti-gaming rules and the attack each answers.
\item What would you tell the desk?
\item In one sentence: what does a router route around?
\end{enumerate}
\end{problem}

\section{Interview questions}

\begin{interviewq}[$\star$ \iqroles{trader}]\label{iq:mx:smart-order-routing:1}
Your sweep shows 5\,000 shares across five venues and fills 2\,600. What happened?
\end{interviewq}

\begin{interviewq}[$\star\star$ \iqroles{developer}]\label{iq:mx:smart-order-routing:2}
How would you implement arrival-time synchronisation in a router, and where does it fail?
\end{interviewq}

\begin{interviewq}[$\star\star$ \iqroles{researcher}]\label{iq:mx:smart-order-routing:3}
How would you estimate a venue's fill probability for passive orders, and why is it hard?
\end{interviewq}

\begin{interviewq}[$\star\star$ \iqroles{researcher}]\label{iq:mx:smart-order-routing:4}
Formulate the allocation of a passive order across venues as an optimisation problem.
\end{interviewq}

\begin{interviewq}[$\star\star$ \iqroles{bank}]\label{iq:mx:smart-order-routing:5}
A client asks whether your router sends its \omterm{def:mx:the-limit-order-book:orders}{limit orders} to the venues that pay you the largest rebates. How do you answer, and how do you show it?
\end{interviewq}

\begin{interviewq}[$\star\star\star$ \iqroles{developer}]\label{iq:mx:smart-order-routing:6}
Design the venue-statistics store of a router: what it records, how it ages the data, and how the router uses it.
\end{interviewq}
